Le collège est l'endroit où le calcul devient démonstration. En sixième, on applique des
règles ; en troisième, on justifie. Les sept chapitres qui suivent parcourent ce
déplacement, du critère de divisibilité par deux jusqu'à la réciproque du théorème de
Pythagore.

Deux choses y frappent quand on les formalise. D'abord, la quantité d'énoncés qui tiennent
en une ligne de bibliothèque parce qu'ils sont vrais dans n'importe quel anneau : la règle
des signes, la distributivité, le produit nul. Ensuite, à l'inverse, le coût réel de ce
qu'on présente comme des évidences — l'unicité de la décomposition en facteurs premiers,
admise au collège, demande le lemme d'Euclide, qui demande le théorème de Gauss.

Les aires et les volumes forment le point dur de cette partie, et pour une raison de fond :
au collège, l'aire du disque \emph{est} une formule qu'on donne ; en théorie de la mesure,
c'est un théorème d'intégration. Le chapitre concerné le dit à sa place.
